\SourceLeaf{35}

Der Wert des Bruchs $\dfrac{P_\theta}{P_{\theta+1}}$ fällt folglich zu klein aus,
wenn man im Zähler die ersten beiden Glieder und im Nenner nur das erste
berücksichtigt. Mit anderen Worten, der wahre Wert dieses Bruchs ist größer als
$(2\theta+1) . \dfrac{n}{mm}\left\{1-\dfrac{1}{2(2\theta+1)}\dfrac{mm}{nn}\right\}$ oder
$\dfrac{n}{mm}(2\theta+1)-\dfrac{1}{2n}$. 
Hieraus folgt, dass $P_\theta, P_{\theta+1}, P_{\theta+2}$ und\,so\,weiter eine stets
abnehmende unendliche Folge von Größen bilden, sofern $2\theta+1$ mindestens
$\dfrac{mm}{n}+\dfrac{mm}{2nn}$ ist. Es ist also unmöglich, dass sämtliche Glieder der Folge
$P, P_1, P_2, P_3$ und\,so\,weiter zu ganzen Zahlen proportional sind. Dies wäre jedoch
notwendig der Fall, wenn $m$ und $n$ ganze Zahlen sind und
die ersten beiden Glieder dieser Folge $P, P_1$ in einem rationalen
Verhältnis zueinander stehen. Es gilt nämlich

\[
\begin{aligned}
P_2 &= nP_1 - mP\\
P_3 &= 3nP_2 - mmP_1\\
P_4 &= 5nP_3 - mmP_2\\
P_5 &= 7nP_4 - mmP_3\\
&\text{und\,so\,weiter,}
\end{aligned}
\]
\noindent allgemein, von der zweiten Formel an,

\[
P_{\theta+2} = (2\theta+1)nP_{\theta+1} - mmP_\theta.
\]

In diesen Formeln kann man die Größen $P, P_1, P_2$ und\,so\,weiter offenbar durch
zu ihnen proportionale Größen ersetzen. Wären aber die ersten beiden davon ganze Zahlen,
so wären auch alle folgenden ganze Zahlen, was, wie eben gezeigt, unmöglich ist.

Sind also $m, n$ ganze Zahlen, so kann der Wert des Bruchs
$\dfrac{P_1}{P}$, also die Tangente des Bogens
$\dfrac{m}{n}$, nicht rational sein.
